\section{Limitations}
\label{sec:limitations}

In Sections~\ref{sec:prompt_and_fewshot} and ~\ref{sec:rai_evals}, we carried out extensive evaluation of all released models at varying scales. We saw parity in performance for standard evaluation datasets used in the GPT-3 models. 
Moreover, we performed safety, bias, and inclusion evaluations, again seeing largely comparable performance with some variations in toxicity and hate speech detection. 
However, such evaluations may not fully characterize the complete limitations of these models. In general, we qualitatively observe that \OPT{} suffers from the same limitations noted in other LLMs \cite{brown2020gpt3,J1WhitePaper,thoppilan2022lamda,gopher2022,megatron2022,palm2022,bender2021dangers}.

In particular, we found \OPT{} does not work well with declarative instructions or point-blank interrogatives. Prompting with such instructions tends to produce a simulation of a dialogue beginning with such an instruction, rather than an execution of the instruction. Future work into instruction learning, in the vein of InstructGPT \cite{ouyang2022training}, may alleviate these limitations.

\OPT{} also tends to be repetitive and can easily get stuck in a loop. While sampling can reduce the incidence rate of repetitive behavior \cite{holtzman2020nucleus}, we anecdotally found it did not eliminate it entirely when only one generation is sampled. Future work may wish to incorporate more modern strategies for reducing repetition and improving diversity, such as unlikelihood training \cite{welleck2019neuraltext} or best-first decoding \cite{meister-etal-2020-best}.

Similar to other LLMs, \OPT{} can produce factually incorrect statements \cite{adiwardana2020meena,brown2020gpt3, roller-etal-2021-recipes,gopher2022,palm2022,thoppilan2022lamda}. This can be particularly harmful in applications where information accuracy is critical, such as healthcare and scientific discovery \cite{weidinger2021ethical}. Recently, several efforts have reported that retrieval-augmented models can improve factual correctness of LLMs \cite{lewis2020retrieval,komeili2021woi,thoppilan2022lamda,borgeaud2021improving,shuster2022language,nakano2021webgpt}. We believe \OPT{} will also benefit from retrieval-augmentation in future iterations.

As shown in Section~\ref{sec:rai_evals}, we also find \OPT{} has a high propensity to generate toxic language and reinforce harmful stereotypes, even when provided with a relatively innocuous prompt \cite{gehman2020realtoxicityprompts}, and adversarial prompts are trivial to find \cite{dinan2021anticipating}. There has been a great deal of work on mitigations for toxicity and biases \cite{dathathri2019plug,dinan2019build,sheng2019woman,dinan-etal-2020-queens,liu2019does,krause2020gedi,xu2020recipes,liang2021towards,dinan2021anticipating,xu-etal-2021-bot,dhamala2021bold,schick2021self,ouyang2022training}. Depending on downstream applications, future uses of \OPT{} may need to employ these or novel mitigation approaches, especially before any real world deployment. Given our primary goal as a replication of \biggpt{}, we choose not to apply these mitigations in this first release.


In summary, we still believe this technology is premature for commercial deployment. Despite including data sheets and model cards, we believe more scrutiny should be afforded to the training data with additional data characterization and selection criteria in order to use data responsibly. The current practice is to feed the model with as much data as possible and minimal selection within these datasets. Despite having comprehensive evaluations, we would ideally have more streamlined and consistent evaluation setups to ensure replicability and reproducibility of evaluation scenarios. Differences in prompting styles and number of shots for in-context learning could create variations that lead to different results. We hope that the public release of the OPT models will enable many more researchers to work on these important issues.
